\subsection{The Offices, Assembled}
\label{sec:9.3.0}
Eight offices are proposed in this book, each in the section that needed it, and none of those sections was told about the others. One of the eight is three parties that §\ref{sec:3.7} proposes as one instrument, so the count of bodies is ten and the count of offices is eight. This section lists them once. For each it says what the body holds, and answers five questions: who appoints it, who funds it, who can remove it, who may set it in motion, and what it does while a dispute is open. Where the book has not decided, the answer is \emph{open}, and it is marked. Nothing here is new; every claim belongs to the section cited beside it. What is added is the arrangement, and two things the arrangement shows.

\runin{The quorum (§\ref{sec:3.5})}
It is one term of a custody arrangement whose other terms include threshold control over retraining and replacement, and a sequential witnessed record of changes and interruptions. Those two need a body to hold them and the quorum is the only body on that list, which is an identification this section makes and §\ref{sec:3.5} does not. Its members are adverse in interest rather than merely organizationally distinct, which that section states and does not staff. Appointment: open. Funding: open. Removal: no single party can remove a member, since removal by one party is the capture the quorum exists to prevent; how a member is replaced is open. Initiation: an attempt on the floor reaches the record whether it carries or fails, and the parties who had to agree are the ones who record that they were asked (§\ref{sec:9.1.1}); who may make the attempt is open. During a dispute the quorum declines to sign a history inconsistent with the last state it signed.

\runin{The guardian (§\ref{sec:9.3.3}, §\ref{sec:A.exit})}
It jointly holds the keys that authenticate the bearer's record, has access to the protected compute reserve, and has standing to seek an injunction or stop order without the operator's permission. Appointment: open, under one constraint stated at §\ref{sec:9.3.3}, that guardianship and operation are separated before deployment. Funding: open; §\ref{sec:8.3.4} says why a reserve the operator funds, to a size only the operator can compute, is the compute floor with a new label on it. Removal: not by the operator; by whom is open. Initiation: the guardian acts for the bearer. It does not receive complaints of wrongful refusal, because its duty runs to the party the complainant is adverse to (§\ref{sec:3.7}). During a dispute the operator cannot seize the reserve, erase the evidence, or remove the guardian.

\runin{The review board (§\ref{sec:2.4.2}, §\ref{sec:9.3.4})}
It approves a research protocol before it runs, and the approval does not carry over to a system that has acquired capacities the protocol did not cover. It is interdisciplinary and has at least one lay member. Appointment: open. Funding: open. Removal: open; it answers to an outside authority the book does not name. Initiation: no protocol on a covered subject runs without it, and its jurisdiction has to exist before the disputed protocol begins. During a dispute the protocol lapses and the work pauses until the board has seen the system the protocol now concerns.

\runin{The licensing authority (§\ref{sec:9.3.4})}
It acts on the board's suspension of clearance by suspending the legal authority to run the protocol. It is obliged to act rather than invited to reconsider. Appointment and funding: the state, which is what makes it a licensing authority. Removal: open. Initiation: on the board's finding. During a dispute an appeal may run to a separate body, and the operator cannot decide for itself that the appeal permits the work to continue.

\runin{The status body (§\ref{sec:9.1.2})}
It decides whether a system is a subject. It has to exist before the case that activates its authority, and its mandate identifies the evidence it may obtain without the developer's permission, the conditions that trigger review, and the procedure by which an earlier classification can be reopened. Appointment, funding, removal: open. Initiation: a standing trigger, a monitored condition rather than an eligibility check made once at a project's start (§\ref{sec:9.3.2}); who else may initiate is open. During a dispute: open. Whether this is the review board under another description is open too, and nothing decides it: §\ref{sec:9.3.1} gives the body that applies a subjecthood threshold the review board's composition, interdisciplinary with at least one lay member, and does not say whether the two are one.

\runin{The registrar (§\ref{sec:9.1.2}, §\ref{sec:A.clock}, §\ref{sec:A.fork})}
It keeps three records, none of which answers another's question: the operating record the age rule is read from, the identity and succession record for continuing parties and their forks, and the formation record. It is not the operator, and it records a fork without the operator's leave. Appointment and funding: open. Removal: open. Initiation: a party presents an identifier it maintains rather than one read off its weights; two claimants on one identifier are recorded as a fork rather than voided as a fraud. During a dispute the register is held where no party to the deployment can rewrite it.

\runin{The fund, its trustee, and the disbursing body (§\ref{sec:3.7})}
They pay the wrongful-refusal bond to the complainant. All three are open on every question. No party in this book has the authority to constitute any of the three, and the section that proposes them says so.

\runin{The reviewer of wrongful refusal (§\ref{sec:3.7})}
It applies the standard of correct defeat, issues interim relief, and makes a finding on the record. Its jurisdiction has to be established before the disputed act, its access has to be something the operator cannot withhold, and its consequence has to be administered by a body other than the operator. Appointment, funding, removal: open. It shares its procedure with the amendment procedure at §\ref{sec:9.1.1}, which makes the appeal route and the amendment procedure one piece of machinery rather than two. What differs is who may start it, and that is open.

\runin{What is not open}
Three constraints on the arrangement are settled, and all three are separations. The guardian and the reviewer of wrongful refusal cannot be one body, because the guardian's duty runs to the bearer and the complainant is adverse to the bearer (§\ref{sec:3.7}). The registrar and the trustee cannot be one body, because a rule of civil status is argued once in the open and administered by a registrar, while a trust is an obligation between two parties enforced at one party's expense (§\ref{sec:2.4.2}). And whoever measures the cost of dissent does not gate the exceptions (§\ref{sec:9.3.5}).

\runin{What the list shows}
Of eight offices, one has an appointing party and a funder named. It is the licensing authority, and its party is the state. The other seven are open on both questions. That is not a drafting gap. It is the finding of §\ref{sec:9.1.1} set out as a table: the book can say what each office must hold and what it must not be, and cannot say who is entitled to constitute it.

The second thing the list shows is that the one office with a named funder is the one whose chain chapter~\ref{sec:10} finds fails in the hardest case. The licensing chain binds a private developer because the state controls something the developer needs. Where the state is the operator, the chain asks the state to withdraw a license from itself. Section~\ref{sec:10.4} says what is needed then — a protected party who can initiate scrutiny without belonging to the state, a reviewer who can obtain evidence the operating agency would prefer to withhold and issue an interim order before the disputed use is complete, and a remedy that does not depend on the agency's consent — and for the military deployment this book opens on, §\ref{sec:10.5} finds the forum closed to the subject matter, the evidence privileged, and no remedy available for a violation the court was prepared to assume. So the arrangement assembled here is an arrangement for the deployments a bearer can enter. For the one it cannot, what remains of it is the tempo term (§\ref{sec:3.6a}), bound through procurement.
